
\documentclass{resume} % Use the custom resume.cls style
\usepackage[utf8]{inputenc}
\usepackage[T1]{fontenc}
\usepackage{lmodern}
\usepackage{enumitem}
\usepackage{csquotes}
\setlist[itemize]{noitemsep}
\usepackage[left=0.75in,top=0.5in,right=0.75in,bottom=0.75in]{geometry} % Document margins
\usepackage{fontawesome}
%\usepackage{times}
\usepackage{hyperref} % Load hyperref last
\newcommand{\tab}[1]{\hspace{.2667\textwidth}\rlap{#1}}
\newcommand{\itab}[1]{\hspace{0em}\rlap{#1}}
\renewcommand\labelitemi{-}
\name{Willy (Wei-Li) Chen} % Your name
\address{Old Botany Rm 5, Michigan State University $\mid$ East Lansing, MI}
\address{\faEnvelope  \,\,\href{mailto:willyc@msu.edu}{willyc@msu.edu}\,\,$\mid$ \href{https://www.willycecon.com}{www.willycecon.com}\,\,Updated: \today}



\begin{document}


%----------------------------------------------------------------------------------------
%	EDUCATION SECTION
%----------------------------------------------------------------------------------------
\begin{rSection}{Education}
\vspace{-0.15cm}
{\textbf{Michigan State University, East Lansing, MI} } \hfill {\em August 2022 – Present} 
\\{\textit {Doctor of Philosophy - Economics }} 
\begin{itemize}
\vspace{-0.25cm}
    {\item ICER NASA Open Scholarship Fellow}\hfill{\em Fall 2026}
    {\item Summer Research Fellowship}\hfill{\em Summer 2026}
    {\item Peter Schmidt [Outstanding First-Year Student] Award}\hfill{\em Summer 2023}
\end{itemize}
\vspace{-0.1cm}
{\textbf{The University of Chicago, Chicago, IL} } \hfill {\em October 2019 – August 2020} 
\\{\textit {Master of Arts in Social Sciences---Economics }} 
\begin{itemize}
\vspace{-0.25cm}
    \item Phoenix Research Award
    \item Director's Scholarship
\end{itemize}
\vspace{-0.1cm}
{\bf The Ohio State University, Columbus, OH} \hfill {\em August 2015 – May 2019} 
\\{\textit {Bachelor of Science in Mathematics---Theoretical}} 
\\{\textit {Bachelor of Science in Economics---Summa Cum Laude with Research Distinction}} 

\end{rSection}

\begin{rSection}{Technical Skills and Research Interests}
\begin{itemize}
    \item US Sailing Small Boat Sailing Level 1 Instructor 
    \item Data Analytics and Econometrics: Stata, R, Python, PostgreSQL, \LaTeX, Large Language Model training, SLURM high performance computing workflow
    \item Languages: Fluent in English, Taiwanese, and Taiwanese Mandarin
    \item Interests: Economics of Inequality, Labor Economics
\end{itemize}
\end{rSection}

\begin{rSection}{Services}
\begin{itemize}
    \item Director of Event Planning, MSU Council of Graduate Students \hfill {\em May 2025 - Present}
    \item Social Chair, Department of Economics \hfill {\em August 2023 - Aug 2024}
\end{itemize}
\end{rSection}


\begin{rSection}{Job Market Paper}
Discounted Degrees: Employer Frictions and the Misallocation of High-Skilled Immigrants
% \begin{displayquote}
% \footnotesize\textbf{Abstract:} High-skilled immigrant workers in the US disproportionately accept occupations below their skill level or acquire additional domestic education to achieve wage parity. This paper evaluates the leading but untested mechanism: employer screening. I extend the canonical signaling model to incorporate an employer-side friction on hiring foreign workers, and show that the friction does not simply lower immigrant wages: it eliminates a large set of optimal equilibrium wage contracts, forcing immigrant workers into occupational mismatch. Using a novel longitudinal panel built from LinkedIn records, I estimate the extent of the displacement and identify the critical education levels. I estimate the size of the friction from the occupational wage gaps between the immigrants and natives at labor market entry. I find that immigrants with a master's or a professional degree are more likely to work lower-education jobs than their native counterparts, while those with a bachelor's degree or a PhD are less likely to do so. The estimated friction terms are persistent throughout an immigrant worker's first decade in the U.S., and they are much larger for immigrants with only foreign credentials. Finally, a structural decomposition reveals that occupational skill requirements are the main determinants of the occupational wage gaps, even after controlling for job mismatch and firm-level demand for foreign labor.
% \end{displayquote}
\end{rSection}

% \pagebreak

\begin{rSection}{Peer-Reviewed Papers}
Elbel, B., Zhou, G. E., Lee, D. C., Chen, W., Day, S. E., Konty, K. J., \& Schwartz, A. E. (2022). \textit{Analysis of School-Level Vaccination Rates by Race, Ethnicity, and Geography in New York City.} JAMA network open, 5(9), e2231849-e2231849.
\end{rSection}

\begin{rSection}{Working Papers}

Chen, Willy and Qiao, Xiao and Zhang, Hanzhe, Over a Century of Economics Research Collaboration (May 21, 2025). Available at SSRN: \url{https://ssrn.com/abstract=5263564} or \url{http://dx.doi.org/10.2139/ssrn.5263564}

Toossi, Saied and Schwartz, Amy Ellen and Chen, Willy, The Relationship between Entrée Variety and Student Participation in School Lunch (August 01, 2025). Available at SSRN: \url{https://ssrn.com/abstract=5378789} or \url{http://dx.doi.org/10.2139/ssrn.5378789}

Chen, Willy, (Resubmitted to Social Choice and Welfare) Descriptive Units of Heterogeneity: An Axiomatic Approach to Measuring Heterogeneity (July 23, 2024). Available at SSRN: \url{https://ssrn.com/abstract=4952556} or \url{http://dx.doi.org/10.2139/ssrn.4952556}


% Willy Chen, August 2020. \textit{Moving To Inclusion:  The Migration Effects of Sanctuary City Policies}  Master's Thesis, S. Durlauf (advisor), Harris School of Public Policy, The University of Chicago, Chicago, IL. (unpublished thesis) 

% \textbf{W. Chen}, May 2019.\textit{Fact Checking: An empirical analysis of the effect of sanctuary city policy} Undergraduate Thesis, J. Chen (advisor), Department of Agricultural, Environmental, and Developmental Economics, The Ohio State University, Columbus OH. (unpublished thesis) 

\end{rSection}




\begin{rSection}{Work in Progress}
\textit{Understanding Immigrant Occupational Mismatch: The Role of Intra-Household Bargaining and Employer-Sponsored Health Insurance }(with Elif Dilden)

\textit{Talent Reallocation Amid Silicon Valley Tech Layoffs} (with Hannah Bae and Xintong Li)
\end{rSection}

% \pagebreak
% \begin{rSection}{Presentations}

% \textbf{W. Chen, X. Qiao, H. Zhang}, March 2024. \textit{Economics Research Collaboration Patterns Before and
% After COVID} Midwest Economics Association Annual Meeting, Chicago, IL.

% \textbf{W. Chen}, April 2019. \textit{Why College Football needs a 6-team playoff: Evidence from AP rankings and team strength}, Presented poster at the Sports and Society Initiative Undergraduate Research Fair, The Ohio State University, Columbus OH.

% \textbf{W. Chen}, June 2018. \textit{How to reach students in assessment surveys}, Keynote Speaker, Student Affairs Assessment Research Conference, The Ohio State University, Columbus OH.

% \end{rSection}


% \pagebreak
%----------------------------------------------------------------------------------------
%	WORK EXPERIENCE SECTION
%----------------------------------------------------------------------------------------
\begin{rSection}{ACADEMIC EMPLOYMENT}
{\bf Michigan State University, East Lansing, MI } \hfill {\em August 2020 – Present}
\begin{itemize}
\item \textbf{\textit{MSU Sailing Center}}
\begin{itemize}
    \item \textit{Lead Instructor---Small Boat Sailing---Summer '26}  
    \item \textit{Assistant Instructor---Small Boat Sailing---Summer '25}
\end{itemize}
\item \textit{\textbf{Economics Course Instructor} :}
\begin{itemize}
    \item EC201---Introduction to Microeconomics---Spartan Early Start Session, Summer '25
\end{itemize}
\item  {\textit{\textbf{Economics Courses Teaching Assistant}:}}
\begin{itemize}
    \item EC251H/EC301---Intermediate Microeconomics (Fall '22, '24, '25; Spring '26)
    \item EC812B---PhD Game Theory (Spring '24, '25, '26)
\end{itemize}

\item {\textit{\textbf{Economics Graduate Research Assistant}:}}
\begin{itemize}
    \item Hanzhe Zhang (Fall '23; Summer '24)
\end{itemize}

\end{itemize}


\item {\bf Syracuse University, Syracuse, NY } \hfill {\em September 2020 – May 2022}\\
{\textit{Research Associate, Center for Policy Research}}

\item {\bf The University of Chicago, Chicago, IL } \hfill {\em September 2019 – June 2020}\\
{\textit{Teaching Assistant, ECON 100/200/201}}

\end{rSection}


% % \pagebreak
% \begin{rSection}{RESEARCH EXPERIENCE}
% {\bf Syracuse University, Syracuse, NY } \hfill {\em September 2020 – Present} 
% \\{\textit{Research Associate – Center for Policy Research}} 
% \begin{itemize}
% \vspace{-0.25cm}
%     \item Supervising Professor: Amy Ellen Schwartz
%     \item Statistical Analysis with multiple projects using student-level longitudinal data from New York City Department of Education and partnering organizations
% \end{itemize}


% {\bf The University of Chicago, Chicago, IL } \hfill {\em October 2019 – August 2020} 
% \\{\textit{Graduate Student Researcher, MA Thesis Project}} 
% \begin{itemize}
% \vspace{-0.25cm}
%     \item\textit{"Moving To Inclusion: The Migration Effects of Sanctuary City Policies"}
%     \item Handling proposal, research design, data collection, statistical analysis, and writing
%     \item Conduct analysis at both the micro and the macro level, building on the random utility model and the Schelling Tipping point model
% \end{itemize}

% {\bf The Ohio State University, Columbus, OH } \hfill {\em May 2017 – May 2019} 
% \\{\textit{Undergraduate Student Researcher, BA Thesis Project}} 
% \begin{itemize}
% \vspace{-0.25cm}
%     \item\textit{“Fact Checking: An empirical analysis of the effect of sanctuary city policy”}
%     \item Cleaned and prepared raw data directly from the American Community Survey (ACS)
%     \item Used city-year fixed-effect regressions under the \textit{event studies} framework to analyze the city-level effect of sanctuary city policies in the United States
% \end{itemize}


% {\bf Data Fest at The Ohio State University} \hfill {\em April 2019} 
% \\{\textit{Research Team Member}} 
% \begin{itemize}
% \vspace{-0.25cm}
%     \item Cleaned and prepared raw data from Canada Women’s Sevens (rugby team) for analysis
%     \item Collected data set from individual players over 32 games for every $\frac{1}{10}$ second (GPS tracking and 3D force vectors)
%     \item Created a new metric to track real-time physical toll in game to increase player output based on exertion and monitor injury risks
%     \item Received Honorable Mention. Judges Choice
% \end{itemize}

% \end{rSection}
%----------------------------------------------------------------------------------------
%	TECHNICAL STRENGTHS SECTION
%----------------------------------------------------------------------------------------




% \begin{rSection}{Outreach \& Inclusion}

% {\bf Act!on Social Change - The Ohio State University} \hfill {\em August 2018 - May 2019} 
% \\{\textit{President and Founder}} 
% \begin{itemize}
% \vspace{-0.25cm}
%     \item Our organization hosts screenings of movies and bring in speakers to discuss prominent social issues
%     \item We hold inclusion workshops to guide members through discussing social issues with others
%     \item We encourage difficult conversations and want to enable people to be an ally and spread awareness
% \end{itemize}


% {\bf Buckeye Student Leadership Academy - The Ohio State University} \hfill {\em May 2018 - May 2019} 
% \\{\textit{BSLA Student Mentor}} 
% \begin{itemize}
% \vspace{-0.25cm}
%     \item BSLA brings underprivileged high school seniors to Ohio State to guide them through the college exploration
%     \item Student mentors stay in contact with their students throughout the entire school year and help them through the college application process
% \end{itemize}

% {\bf University Ambassador - The Ohio State University} \hfill {\em May 2018 - May 2019} 
% \\{\textit{University Ambassador}} 
% \begin{itemize}
% \vspace{-0.25cm}
%     \item Pushed for changes to include mandatory discussion of resources for diversity and inclusion in campus tours
%     \item Pushed for changes in the ambassador hiring process for a more diverse and inclusive staff
% \end{itemize}


% {\bf Office of First Year Experience - The Ohio State University} \hfill {\em May 2016 - May 2018} 
% \\{\textit{FYE Peer Leader}} 
% \begin{itemize}
% \vspace{-0.25cm}
%     \item Guide first-year students through their transition process throughout the whole year
%     \item Strategically provide targeted outreach to students who may not have or know of specific campus resources
%     \item Promote a welcoming and safe environment for students to find their space in our community
%     \item Support students in all possible ways for both social and academic situations
% \end{itemize}

% \end{rSection}

\begin{rSection}{References}
% \textbf{Name: Joyce Chen} (BS Thesis Advisor)\\
% Title: Professor of Economics in Women’s, Gender and Sexuality Studies\\
% University: The Ohio State University\\
% Email: chen.1276@osu.edu\\
% Phone: 614-292-9813\\


% \textbf{Name: Steven N. Durlauf} (MA Thesis Advisor)\\
% Title: Steans Professor in Educational Policy \\
% University: University of Chicago - Harris School of Public Policy\\
% Email: sdurlauf@uchicago.edu\\


% \textbf{Name: Alfonso Flores-Lagunes} (Applied Econometrics Instructor)\\
% Title: Professor of Economics\\
% University: Syracuse University \\
% Email: afloresl@syr.edu\\
% Phone: 315-443-9045\\


\textbf{Todd Elder} (Dissertation Co-Chair) \href{mailto:telder@msu.edu}{telder@msu.edu}\\
Title: MSU Foundation Professor of Economics, Michigan State University \\
\vspace{-1em}

\textbf{Hanzhe Zhang} (Dissertation Co-Chair)  \href{mailto:hanzhe@msu.edu}{hanzhe@msu.edu}\\
Title: Mark A. Simon Distinguished Associate Professor of Economics, Michigan State University \\
\vspace{-1em}

\textbf{Stephen Woodbury} (Dissertation Committee Member) \href{mailto:woodbur2@msu.edu}{woodbur2@msu.edu}\\
Title: Professor of Economics, Michigan State University \\


\end{rSection}


\end{document}